\chapter{Implementation and Results} 
\label{chap:implementation_results} 
% TODO: Confirm whether this chapter is included in a master document; no complete master file was found in the workspace.
This chapter applies the methodology of Chapter \ref{chap:methodology} to a physical \gls{acr:SUT}. It documents the implemented environment, the instantiation of the method, the collected evidence, the verification of the derived requirements, a discussion of the resulting findings, the validation of the methodology against the thesis objective, and the reliability of the findings. 

The practical work covers the direct-access variant \gls{acr:NV1} defined in Section \ref{subsec:nv1}. The passive-observation variant \gls{acr:NV2} and the in-path variant \gls{acr:NV3} are part of the methodology but were not realised in the laboratory. The validation therefore covers the core \gls{acr:NV1} process and those requirements whose acceptance conditions are observable from the direct-access position; the tests that depend on \gls{acr:NV2} or \gls{acr:NV3} remain specified but not executed.



\section{Scope of the Practical Demonstration}
\label{sec:testbed_setup} 

This section defines the object, the environment, and the boundaries of the practical demonstration. Section \ref{subsec:test_object} identifies the \gls{acr:SUT} and its externally accessible attack surface. Section \ref{subsec:realized_environment} describes the physically implemented test topology. Section \ref{subsec:practical_scope} states which parts of the methodology were realised in the laboratory and which remain specified but not executed.

\subsection{System Under Test} 
\label{subsec:test_object} 

The \gls{acr:SUT} is the \gls{acr:SRIO}, a fieldbus-based safety component that exchanges safety-related input and output data with a higher-level safety controller over PROFINET and PROFIsafe. The evaluated unit identified itself over the \gls{acr:IoT}-Core interface as product \texttt{ifm-AL400S}, hardware version 1.4.0.2, hardware revision AA, and software revision 1.1.4.888. All reported results refer exclusively to this configuration.

Following the asset classification of Section \ref{subsec:assets_attack_surfaces}, the \gls{acr:SRIO} is decomposed into hardware, software, and data assets. For the practical assessment, the externally accessible attack surface was divided into the \gls{acr:HTTP}-based \gls{acr:IoT}-Core diagnostic interface, the \gls{acr:PROFINET} and PROFIsafe communication path, the \gls{acr:DCP} configuration functions, and the firmware update mechanism. Only the interfaces accessible from the implemented \gls{acr:NV1} position were exercised. Internal debug and inter-processor interfaces, such as the \gls{acr:SysCom} link between the \gls{acr:COM} and \gls{acr:SCPU} modules and the \gls{acr:JTAG} interface, are included in the asset model but were not accessible within the practical test scope. %TODO warum waren sie nicht anwendbar? z.B. JTAG, physischer Zugriff --> weil SRIO umspritzt
%TODO hab ich die Assets hier richtig definiert nach Grundlagen subsec:assets_attack_surfaces}
%TODO sind das alles IT/OT Schnittstellen?

Table \ref{tab:assets} summarises the assets and their protection objectives. It is not exhaustive, but it covers the main elements that are relevant for the security assessment.
%TODO überprüfen ob das Asset Mapping zur Theorie passt: subsec:assets_attack_surface
{\renewcommand{\arraystretch}{1.2}
	\small 
	\begin{xltabular}{\textwidth}{@{} 
			c 
			>{\raggedright\arraybackslash\hsize=0.9\hsize}X 
			>{\raggedright\arraybackslash\hsize=1.4\hsize}X 
			>{\raggedright\arraybackslash\hsize=0.7\hsize}X 
			@{}}
		
		% --- Kopf- und Fußzeilen-Definition für die Longtable ---
		
		% 1. Kopfzeile auf der ersten Seite (inkl. Caption)
		\caption{Overview of assets and protection objectives}
		\label{tab:assets} \\
		\toprule
		\textbf{ID} & \textbf{Asset} & \textbf{Description} & \textbf{Protection \newline objectives} \\ 
		\midrule
		\endfirsthead
		
		% 2. Kopfzeile auf allen folgenden Seiten
		\caption*{Overview of assets and protection objectives (continued)} \\
		\toprule
		\textbf{ID} & \textbf{Asset} & \textbf{Description} & \textbf{Protection \newline objectives} \\ 
		\midrule
		\endhead
		
		% 3. Fußzeile auf allen Seiten außer der letzten
		\midrule
		\multicolumn{4}{r}{\footnotesize\textit{Continued on next page}} \\
		\endfoot
		
		% 4. Fußzeile auf der allerletzten Seite der Tabelle
		\bottomrule
		\endlastfoot
		
		% --- Eigentlicher Tabelleninhalt ---
		
		A & Trusted Safety Function & \gls{acr:SRIO} correctly executes its specified safety function. & Integrity, \newline Availability \\ 
		
		B & Integrity of Safety Configuration & Safe behaviour is determined by the intended configuration. & Accountability, \newline Authorisation, \newline Integrity, \newline Authenticity \\ 
		
		C & Integrity and Authenticity of Safety-Relevant Process Data & Input, output and PROFIsafe data correspond to the actual safety state. & Integrity, \newline Authenticity, \newline Availability \\ 
		
		D & Authenticity and Integrity of Safety Software & Bootloader and firmware remain authentic and unmodified. & Accountability, \newline Authenticity, \newline Integrity \\ 
		
		E & Integrity of Safety Monitoring & Self-tests, plausibility checks, diagnostics and fault responses remain trustworthy. & Integrity, \newline Availability \\ 
		
		F & Separation of Safety and Non-Safety Domain & A compromised \gls{acr:COM} system must not affect the integrity of the safety function. & Confidentiality, \newline Integrity, \newline Availability \\ 
		
		G & Integrity of Operating Mode & Test and update functions must not be activated or used without authorisation. & Accountability, \newline Authorisation, \newline Integrity, \newline Authenticity, \newline Availability \\ 
		
		H & \gls{acr:SRIO} Functionality & Availability of \gls{acr:SRIO} functionality shall be ensured. & Availability \\ 
		
		I & Integrity, Availability and Confidentiality of Audit/Tracing Data & Evidence of interventions, configuration changes and installed software versions must be generated, retained for the mandated period and access-restricted. & Accountability, \newline Availability, \newline Confidentiality \\
		
	\end{xltabular}
}


\subsection{Implemented Test Environment}
\label{subsec:realized_environment}
The test environment used one switch and one flat address range. Figure \ref{fig:testbed_topology} shows the physical topology, and Table \ref{tab:testbed_components} lists the components and their roles.

\begin{table}[H]
	\centering
	\small
	\caption{Components of the implemented test environment}
	\label{tab:testbed_components}
	\begin{tabularx}{\textwidth}{l l c X}
		\toprule
		\textbf{Component} & \textbf{Address} & \textbf{Port} & \textbf{Role} \\
		\midrule
		Industrial switch & 192.168.0.91 & - & Layer-2 connectivity for all nodes \\
		\gls{acr:PLC} (F-Host) & 192.168.0.1 & 1 & Safety controller and PROFINET/PROFIsafe counterpart \\
		\gls{acr:SRIO} (\gls{acr:SUT}) & 192.168.0.2 & 2 & Device under test \\
		Kali Linux test system & 192.168.0.80 & 3 & Security test system, connected through a USB 3.0 to Gigabit Ethernet adapter \\
		Company laptop & 192.168.0.7 & 4 & Engineering station providing \acrshort{acr:TIA} Portal for planned identification and maintenance read operations \\
		\bottomrule
	\end{tabularx}
\end{table}

All nodes shared the 192.168.0.0/24 segment and one Ethernet broadcast domain. No \gls{acr:VLAN} separation was configured. The test system reached the \gls{acr:SUT} and the controller directly over \gls{acr:IP}. The test object was the \gls{acr:SRIO} on port 2. The test initiator was the Kali system on port 3. The company laptop on port 4 served as the engineering station and provided \gls{acr:TIA} Portal for planned identification and maintenance read operations. These read operations were not part of the evidence-backed execution subset reported in this chapter. The available observation points were the direct request and response traffic of the test system and the \gls{acr:HTTP} responses of the \gls{acr:IoT}-Core service. Traffic that was not addressed to the test system was not necessarily visible, because a switch forwards unicast frames only to the destination port. \gls{acr:NV1} was therefore the realised position. \gls{acr:NV2} and \gls{acr:NV3} were not implemented.

\begin{figure}[H]
	\centering
	\begin{tikzpicture}[
		% Definition der Design-Stile für die Bachelorarbeit
		node distance=1cm and 1cm,
		box/.style={
			draw=black!70, 
			thick, 
			rectangle, 
			rounded corners=1.5mm, 
			minimum width=4.2cm, 
			minimum height=1.8cm, 
			align=center, 
			drop shadow={opacity=0.15}, 
			fill=white,
			font=\small
		},
		switch/.style={
			box, 
			fill=gray!10, 
			minimum height=2cm, 
			minimum width=4.5cm
		},
		link/.style={
			draw=black!80, 
			thick, 
			<->, % Bidirektionale Netzwerkverbindung
			>=Latex
		},
		port/.style={
			fill=white, 
			inner sep=2pt, 
			font=\scriptsize\bfseries,
			text=black!80
		}
		]
		
		% Zentrales Element: Der Switch
		\node[switch] (switch) {
			\textbf{Industrial Switch}\\
			192.168.0.91\\[1.5mm]
			\textit{\scriptsize Layer-2 connectivity}
		};
		
		% Periphere Knoten (Nodes)
		\node[box, above left=of switch] (plc) {
			\textbf{PLC (F-Host)}\\
			192.168.0.1\\[1.5mm]
			\textit{\scriptsize Safety controller}
		};
		
		\node[box, above right=of switch] (sut) {
			\textbf{SRIO (SUT)}\\
			192.168.0.2\\[1.5mm]
			\textit{\scriptsize Device under test}
		};
		
		\node[box, below left=of switch] (kali) {
			\textbf{Kali Linux Test System}\\
			192.168.0.80\\[1.5mm]
			\textit{\scriptsize Security test system}
		};
		
		\node[box, below right=of switch] (laptop) {
			\textbf{Company Laptop}\\
			192.168.0.7\\[1.5mm]
			\textit{\scriptsize Engineering station}
		};
		
		% Verbindungen (Kabel) ziehen und Ports am Switch beschriften
		% pos=0.75 platziert das Port-Label nahe am Switch
		\draw[link] (plc.south east) -- (switch.north west) 
		node[pos=0.75, port, xshift=-0.1cm, yshift=0.2cm] {Port 1};
		
		\draw[link] (sut.south west) -- (switch.north east) 
		node[pos=0.75, port, xshift=0.1cm, yshift=0.2cm] {Port 2};
		
		\draw[link] (kali.north east) -- (switch.south west) 
		node[pos=0.75, port, xshift=-0.1cm, yshift=-0.2cm] {Port 3};
		
		\draw[link] (laptop.north west) -- (switch.south east) 
		node[pos=0.75, port, xshift=0.1cm, yshift=-0.2cm] {Port 4};
		
	\end{tikzpicture}
	\caption{Logical topology and addressing of the implemented test environment}
	\label{fig:testbed_topology}
\end{figure}

The device documentation defines organisational preconditions for secure operation. The intended use assumes the device in a protected zone behind a firewall. A firmware update is expected to be performed by a trusted person behind that firewall. A security risk assessment is required. The debug interface is primarily a physical interface. These preconditions bound the threat model of the practical tests.
These preconditions are relevant to the compensating-countermeasure allocation: zone protection, trusted update procedures, and the system-level security risk assessment may address obligations that the component does not implement itself. They therefore inform the interpretation of component-capability findings in Section \ref{sec:requirement_verification}.



\subsection{Practical Coverage and Limitations} 
\label{subsec:practical_scope} 

The demonstration applied the methodology to the \gls{acr:SRIO} under the \gls{acr:NV1} network position defined in Section \ref{subsec:nv1}. The test system acted as a direct \gls{acr:IP} participant on the shared switch. The passive-observation variant \gls{acr:NV2} and the in-path variant \gls{acr:NV3} were not realised. No mirror or \gls{acr:SPAN} port was configured, because the available laboratory switch does not support port mirroring, and the test system was never placed inline, because rebuilding the setup for an inline bridge exceeded the available time. The cyclic \gls{acr:PROFINET} and PROFIsafe communication between the controller and the \gls{acr:SUT} was therefore not observed, and no claim is made about that traffic.

The empirical coverage was bounded by the single test object, the implemented topology, the available interfaces, and the \gls{acr:NV1}-compatible activities. Within \gls{acr:NV1}, the evaluation focused on the \gls{acr:IoT}-Core \gls{acr:HTTP} interface and on tests in which the test system acted as an unauthorised communication partner. Where a company tool assumed an interface that the device does not expose, the test was recorded as an applicability finding rather than as a failure, following Section \ref{subsec:company_framework}. The scope of this practical validation must be distinguished from the achievement of the thesis objective, which is to develop and validate a systematic methodology for verifying security functions rather than to execute every technically possible test.



\section{Instantiation of the Methodology} 
\label{sec:instantiation} 

This section instantiates the abstract methodology of Chapter \ref{chap:methodology} on the \gls{acr:SRIO}. Section \ref{subsec:req_asset_interface} binds each of the fourteen derived requirements (Table \ref{tab:mvo-requirements}) to the affected asset, the security function that treats it, the concrete interface, and the network access variant required for its verification. Section \ref{subsec:cm_applicability} applies the applicability assessment defined in Section \ref{subsec:company_framework} to the ten company countermeasure tests. Section \ref{subsec:cm_tc_map} consolidates both mappings into the set of test cases that were actually selected for execution.

\subsection{Requirement, Asset and Interface Mapping}
\label{subsec:req_asset_interface}
%TODO stimmt nicht so ganz: in subsec:asset_threat_mapping hab ich eigentlich nur beschrieben, dass ich assets und threats mappen muss, hier wird es wirklich umgestzt ich habe assets für das srio besimmt, dann threats (stride), gaps und mitigations von dem srio aktuell, abstrakte test scenarien, abstrakte testfälle --> testumgebung gestimmen
The instantiation binds the abstract asset-and-threat mapping of Section \ref{subsec:asset_threat_mapping} to the \gls{acr:SRIO}. Each requirement is paired with the affected asset (Table \ref{tab:assets}), the security function that treats it (Section \ref{subsec:security_functions}), the interface on which it is exercised, and the network access variant that the corresponding test requires. Table \ref{tab:req_asset_interface} shows the complete mapping for all fourteen requirements. The underlying derivation chain from the Regulation through prEN 50742 and \gls{acr:IEC} 62443-4-2 is documented in Appendix \ref{app:requirement_derivation}. 
%TODO stimmt das mapping? workflow.md ; ist das mapping sinnvoll? 
\begin{table}[H] 
	\centering 
	\small 
	\caption{Instantiated requirement-to-asset-to-interface mapping for the \gls{acr:SRIO}} 
	\label{tab:req_asset_interface} 
	\begin{tabularx}{\textwidth}{l l X X l} 
		\toprule 
		\textbf{Req.} & \textbf{Asset} & \textbf{Security function} & \textbf{Interface} & \textbf{Access} \\ 
		\midrule 
		RQ-001 & F & Availability and resilience (domain separation) & Non-safety path (\gls{acr:IoT}-Core) and safety communication path & NV1, NV2 \\ 
		\addlinespace 
		RQ-002 & C & Integrity protection & Firmware update service; PROFINET/PROFIsafe process data & NV1, NV3 \\ 
		\addlinespace 
		RQ-003 & I & Logging and auditing & IoT-Core error log & NV1 \\ 
		\addlinespace
		RQ-004 & D & Integrity protection / identification & IoT-Core \texttt{deviceinfo}/\texttt{firmware} & NV1 \\ 
		\addlinespace 
		RQ-005 & D & Integrity protection & Firmware update service & NV1, NV3 \\ 
		\addlinespace 
		RQ-006 & D & Identification & IoT-Core \texttt{firmware} & NV1 \\ 
		\addlinespace 
		RQ-007 & D, H & Availability & IoT-Core \texttt{deviceinfo}, HTTP service & NV1 \\ 
		\addlinespace RQ-008 & I & Logging and auditing & IoT-Core error log & NV1 \\ 
		\addlinespace 
		RQ-009 & I, G & Logging and auditing; authorisation & IoT-Core error log; configuration interfaces & NV1 \\ 
		\addlinespace 
		RQ-010 & H & Availability and resilience & HTTP service and safety communication & NV1, NV2 \\ 
		\addlinespace 
		RQ-011 & B & Authorisation & Configuration interfaces (DCP, engineering tool) & NV1 \\ 
		\addlinespace 
		RQ-012 & I & Logging and auditing (accountability) & IoT-Core error log & NV1 \\ 
		\addlinespace 
		RQ-013 & D, I & Identification; logging (version history) & IoT-Core firmware version history & NV1 \\ 
		\addlinespace 
		RQ-014 & I & Confidentiality; authorisation & Tracing-log access; HTTP transport & NV1 \\ 
		\bottomrule 
	\end{tabularx} 
\end{table}



\subsection{Applicability Assessment}
\label{subsec:cm_applicability}

Table \ref{tab:cm_applicability} records the applicability and execution status of the ten company tests. Execution status and requirement assessment are separate: a tool may be applicable and executed without producing sufficient evidence for a requirement verdict. Not-applicable outcomes are properties known in advance from the device or topology, not product failures. CM-002 requires a second network zone, which the flat segment does not provide, CM-003 assumes an authentication endpoint that the \gls{acr:SUT} does not expose, and CM-006 assumes an operating-system shell that the \gls{acr:SUT} does not expose. CM-005 was directed at \gls{acr:TCP} port 80 and showed that the tested endpoint did not provide \gls{acr:TLS}. This is an interface observation, not a \gls{acr:TLS} test of \gls{acr:TCP} port 443.

\begin{table}[H] 
	\centering 
	\small 
	\caption{Applicability of the company framework tests to the \gls{acr:SUT}} 
	\label{tab:cm_applicability} 
	\begin{tabularx}{\textwidth}{l l X} \toprule \textbf{CM} & \textbf{Tool} & \textbf{Applicability outcome} \\ 
		\midrule 
		CM-001 & Nmap & Applicable; performed \\ 
		CM-002 & Nmap & Not applicable to the topology; no zone or conduit model in a flat segment \\ 
		CM-003 & Hydra & Not applicable; no authentication endpoint is exposed; anonymous read access was observed separately \\ 
		CM-004 & \gls{acr:OWASP} ZAP & Applicable with adaptation; web scan performed, multi-role comparison not applicable \\ 
		CM-005 & SSLscan & Not applicable to the original procedure; \gls{acr:TCP}/80 provided \gls{acr:HTTP} without a \gls{acr:TLS} handshake or certificate \\ CM-006 & Lynis & Not applicable to the device; no operating-system shell \\ 
		CM-007 & Nikto & Applicable; performed \\ 
		CM-008 & OpenVAS/GVM & Blocked; scanner could not be made operational in the laboratory setup \\ 
		CM-009 & Siege & Applicable; performed, load summary recorded; 50-connection run outside the specified envelope \\ 
		CM-010 & boofuzz & Applicable; performed \\ 
		\bottomrule 
	\end{tabularx} 
\end{table}

\subsection{Selected Test Cases and Execution Status}
\label{subsec:cm_tc_map}
Table \ref{tab:cm_tc_map} maps the company tests to the derived requirements and to the security category in which they are reported. One practical activity can support several requirements. The mapping is not one-to-one. For example, one Nmap run supports the exposure countermeasure CM-001 and the identification and asset-discovery category, and one query of the identification interface supports several identification requirements. The tests CM-002, CM-006, and CM-008 do not appear in the mapping, because they are not applicable or could not be operated on this device and therefore support no requirement. 

The table also lists requirements for which no applicable company test covered the complete acceptance condition under \gls{acr:NV1}. RQ-002 and RQ-005 belong to the firmware and configuration-integrity category (Section \ref{subsec:cat_config_firmware}). Baseline and event-surface reads were performed for RQ-008 and RQ-009, but no complete software-installation or configuration-modification scenario was induced.
%TODO stimmt die Tabelle?, wo ist CM-003?
\begin{table}[H]
	\centering
	\small
	\caption{Mapping of company tests, requirements, and security categories}
	\label{tab:cm_tc_map}
	\begin{tabularx}{\textwidth}{l l X X}
		\toprule
		\textbf{CM} & \textbf{Req.} & \textbf{Test category} & \textbf{Evidence} \\
		\midrule
		CM-001, CM-007 & RQ-001 & Identification and asset discovery & Nmap, Nikto \\
		- & RQ-004, RQ-006, RQ-007, RQ-013 & Identification and asset discovery & \gls{acr:IoT}-Core inventory \\
		CM-005 & - & Communication security & SSLscan; defence-in-depth observation only \\
		- & RQ-003, RQ-008, RQ-012 & Logging and traceability & error log, systick; baseline and intervention evidence \\
		CM-009 & RQ-007, RQ-010 & Availability and resilience & Siege \\
		CM-010 & RQ-010 & Robustness & boofuzz \\
		CM-004 & RQ-011, RQ-014 & Authorisation and access control & ZAP, \gls{acr:IoT}-Core; Hydra not applicable \\
		- & RQ-002, RQ-005 & Firmware and configuration integrity & no complete execution under NV1 \\
		- & RQ-009 & Logging and traceability & baseline only; modification event not induced \\
		\bottomrule
	\end{tabularx}
\end{table}




\section{Test Execution and Results}
\label{sec:executed_tests} 
%TODO auf neue Ergebnisse anpassen
The two-level test catalogue derived in Chapter \ref{chap:methodology} defines abstract test scenarios for all fourteen requirements, including the tests that depend on the passive-observation variant \gls{acr:NV2} and the in-path variant \gls{acr:NV3}. Only a subset of these was executed in the practical phase: the applicable company tests and the requirement-oriented tests that the NV1 position supports. This section groups the tests into the security categories defined in Section \ref{subsec:test_categorisation} and describes them at the level of the test activity, the method, and the result, rather than by individual test-case identifier. Each category summarises the abstract scenarios it contains, and then reports the concrete execution and the stored evidence for the activities performed under \gls{acr:NV1}. The tests that require NV2 or NV3 are described as designed but not executed. Every reported result is limited to what the stored evidence shows. A tool execution is not treated as requirement verification, and an expected outcome from a test specification is not presented as an observed result. 

Each category follows the same logic: it states the machine-level obligation from Regulation (EU) 2023/1230 from which the component requirement is derived, what must therefore be checked, how the property can be tested and which company countermeasure applies, and what the general finding on the device was. Every performed test is stored in some form as evidence, depending on the test as a text file, an HTML report, a log file, or a screenshot. The exact tool commands for the executed activities are listed in Appendix \ref{app:tool_commands}.



\subsection{Identification and asset discovery}
\label{subsec:cat_identification}

This category covers the discovery of the reachable attack surface and the identification of the installed software and data, addressing RQ-001 on the non-safety path and requirements RQ-004, RQ-006, RQ-007, and RQ-013. It targets the separation of the safety and non-safety domains and the software and data-identification assets. The abstract catalogue contains the attack-surface enumeration, the identification-surface inventory, the \gls{acr:SBOM} completeness analysis, the plaintext-transport observation, the data-signature uniqueness check, and the version-history and per-upload-timestamp checks. The bill-of-materials, signature-uniqueness, and version-history scenarios were designed but not executed. The activities executed under NV1 are reported below.

\paragraph{Attack-Surface Enumeration}
\label{subsec:res_attack_surface}
This activity examined the externally reachable services on the non-safety path. It relates to CM-001 and CM-007, to requirement RQ-001, and to asset F, the separation of the safety and the non-safety domain. It targets the \gls{acr:IoT}-Core \gls{acr:HTTP} path and was run under \gls{acr:NV1}. The commands are listed in Appendix \ref{app:cmd_attacksurface}.

A full \gls{acr:TCP} port scan and a top-ports \gls{acr:UDP} scan were run against the \gls{acr:SUT} with Nmap. The \gls{acr:TCP} scan covered all 65535 ports and reported one open port, port 80, which returned \gls{acr:HTTP} responses without a server banner. All other \gls{acr:TCP} ports were reported as closed. The open \gls{acr:HTTP} port is consistent with the \gls{acr:IoT}-Core diagnostic interface. The \gls{acr:UDP} scan of the top 200 ports reported port 161 as open, identified as \gls{acr:SNMP}, and port 49152 as open or filtered. An open \gls{acr:SNMP} port is consistent with a PROFINET device, because the PROFINET standard uses \gls{acr:SNMP} for network diagnostics. The reported open-port count depends on the scan timing. At the slower \texttt{-T3} timing the scan reported the stable result above, whereas the faster \texttt{-T5} timing reported further open ports. The additional ports did not reappear on repetition, which suggests that the faster timing affects the observed responses rather than revealing additional stable services. Packet loss, SYN-backlog exhaustion, rate limiting, device overload, or test-system effects are alternative explanations. If aggressive timing does destabilise the device's responses, that would itself be an availability observation relevant to RQ-010; this remains an open item.

The externally reachable services on the non-safety path are the \gls{acr:HTTP} service on port 80 and the \gls{acr:SNMP} service on port 161. Nikto targets the exposure of debug and service interfaces on a web server, which is the objective of the company countermeasure CM-007. On this device the physical debug interface is encapsulated and not reachable, so the scan characterises only the reachable \gls{acr:HTTP} service surface. The Nikto report lists missing \gls{acr:HTTP} security headers, including content-security-policy, x-content-type-options, strict-transport-security, and permissions-policy. These headers are a defence-in-depth measure, so their absence is an observation rather than a direct vulnerability. The report also flags several \texttt{/actuator/*} paths as Spring Boot Actuator endpoints that return \gls{acr:JSON}.

The Nikto findings are unvalidated scanner results \cite[p. 39]{richier2011security}. The Spring Boot Actuator finding is a probable false positive. The device is an ifm \gls{acr:IoT}-Core and not a Spring Boot application, and its \gls{acr:HTTP} service returns a generic \gls{acr:JSON} structure for many paths, which a generic scanner signature can misclassify as an exposed Actuator endpoint. The finding therefore does not, on its own, indicate a risk. Directly requesting one such path and inspecting the raw response remains an open verification item; no result is claimed here.

The enumeration supports RQ-001 at the level of the non-safety path. It shows which services a connecting device can reach. It does not, on its own, demonstrate a hazardous situation. The \gls{acr:SNMP} service on port 161 was deliberately excluded from the practical scope because it was not examined beyond enumeration. Its authentication, information-disclosure and logging implications remain open items for a future test campaign.

\paragraph{Identification-Surface Inventory}
\label{subsec:res_identification}
This activity enumerated the identification data that the device exposes. It relates to requirements RQ-004, RQ-006, RQ-007, and RQ-013, and to assets B and D. It targets the \gls{acr:IoT}-Core \texttt{deviceinfo} and \texttt{firmware} paths and was run under \gls{acr:NV1}. The commands are listed in Appendix \ref{app:cmd_inventory}.

The identification interface was queried over \gls{acr:HTTP}. The interface required no credentials, so the read operation was anonymous. For each queried path the device returned a \gls{acr:JSON} response with status code 200 that contained the requested identifier. The read operations covered the hardware version and revision, the software revision, the \gls{acr:COM} and \gls{acr:SCPU} firmware and bootloader versions, the host firmware version, and the fieldbus firmware version. A successful read operation is recognisable by the returned status code 200, which indicates a successful response. A separate single read of the software revision also returned status 200, which documents the baseline availability of the interface.

The device exposes top-level version identifiers for the \gls{acr:COM} firmware, the \gls{acr:SCPU} firmware, the bootloaders, and the fieldbus firmware, all through anonymous read access. The inventory supports RQ-004, RQ-006, and RQ-007 at product granularity. It confirms that safety-relevant software is identifiable and available on request. It does not establish a complete \gls{acr:SBOM}. The versions of third-party components, such as the communication stack or the real-time operating system, are not individually exposed. The visible data is therefore an identification surface rather than a complete installed-software inventory or a full \gls{acr:SBOM}.


\subsection{Communication security}
\label{subsec:cat_communication}

This category covers the protection of the communication paths against corruption and disclosure, addressing RQ-001, RQ-002, and the transport aspect of RQ-014. It targets the safety-relevant process data and the audit data. The abstract catalogue contains the transport-confidentiality observation on the non-safety path, the integrity-rejection and in-path manipulation scenarios for the safety telegram, and the watchdog-passivation scenario. The safety-telegram integrity, in-path manipulation, and watchdog scenarios require the \gls{acr:NV2} or \gls{acr:NV3} position or the engineering tool and were designed but not executed. The transport observation executed under NV1 is reported below.

\paragraph{Cleartext Transmission of the Identification Data}
\label{subsec:res_cleartext}
This activity examined the transport protection of the identification data. It relates to CM-005, to requirements RQ-004 and RQ-014, and to asset I. It targets the \gls{acr:IoT}-Core \gls{acr:HTTP} service and was run under \gls{acr:NV1}. The commands are listed in Appendix \ref{app:cmd_cleartext}.

SSLscan verifies the \gls{acr:TLS} configuration of a \gls{acr:TLS} endpoint. It probes for a \gls{acr:TLS} handshake on a given host and port and reports the enabled protocol versions and the certificate. The tool was directed at the \gls{acr:HTTP} service on port 80. The saved output shows that no SSL %TODO stimmt SSL?
or \gls{acr:TLS} protocol version is enabled and that no certificate can be parsed. The output therefore confirms that port 80 does not provide a \gls{acr:TLS} endpoint, so the service is plain \gls{acr:HTTP} without transport encryption. The identification data of Section \ref{subsec:res_identification} observed in the test system's own requests and responses was therefore transmitted in plaintext.
The absence of \gls{acr:TLS} is a device characteristic, not a failed \gls{acr:TLS} test. It is relevant to RQ-014 as a defense-in-depth observation. %TODO was bedeutet defense-in-depth?
A confidentiality impact requires two further conditions. The transmitted data must be sensitive, and an unauthorised party must be able to observe it. The second condition depends on the \gls{acr:NV2} position, which was not realised. The confidentiality claim is therefore limited to the cleartext transmission that the test system observed in its own requests and responses.



\subsection{Logging and traceability}
\label{subsec:cat_logging}
This category covers the evidence, accountability, and retention requirements RQ-003, RQ-008, RQ-009, and RQ-012, with the audit-data asset. The abstract catalogue contains the intervention-evidence baseline, the time-attribution and cold-start-deletion scenarios, the software- and configuration-modification evidence scenarios, and the circular-buffer and retention scenarios. The software-installation and configuration-modification evidence scenarios require an induced software or configuration change and were not fully executed. The intervention-evidence and time-attribution activity executed under \gls{acr:NV1} is reported below; the log-access probe is reported under the authorisation and access-control category.

\paragraph{Intervention Evidence and Time Attribution}
\label{subsec:res_evidence}
This activity examined what evidence the device records for an intervention and how it is time-referenced. It relates to requirements RQ-003, RQ-008, and RQ-012, and to assets I and G. It targets the \gls{acr:IoT}-Core \texttt{errorlog} and \texttt{systick} paths and was run under \gls{acr:NV1}. The commands are listed in Appendix \ref{app:cmd_evidence}.

The machine-level requirement requires a tracing-log capability, and usable intervention evidence requires a time reference that remains unambiguous across the relevant operating states. No dedicated tracing log was exposed through the examined interface. The available substitute artefact is an error log with a relative system tick, and it is cleared at a cold restart. The activity therefore records both the apparent absence of the required tracing-log function and the properties of the substitute artefact.

The error log and the system tick counter were read over \gls{acr:HTTP} in a controlled sequence. The sequence recorded a baseline, an induced input fault by short-circuiting pins 1 and 3 of a digital input of the \gls{acr:SRIO}, a cold restart of the device by removing power, and a second induced input fault. Table \ref{tab:rq003_evidence} lists the recorded system tick values and the number of log entries at each step.

\begin{table}[H]
	\centering
	\small
	\caption{Recorded system tick values and error-log entry counts}
	\label{tab:rq003_evidence}
	\begin{tabular}{l r r}
		\toprule
		\textbf{Step} & \textbf{System tick} & \textbf{Log entries} \\
		\midrule
		Baseline & 20969 & 0 \\
		After event 1 (input fault) & 68629 & 2 \\
		After cold restart & 20459 & 0 \\
		After event 2 (input fault) & 44869 & 2 \\
		\bottomrule
	\end{tabular}
\end{table}


At the baseline the error log was empty. The short-circuit at the digital input triggered a fault, which produced two error entries on the safety \gls{acr:CPU}. Each entry carried a relative system tick value but no absolute timestamp and no actor identity. After the cold restart the error log was empty again, and the system tick had dropped, as Table \ref{tab:rq003_evidence} shows. The second input fault produced two new entries with the same recorded error code. Because the system tick is reset by a restart, an entry recorded after the restart can carry a lower tick than an earlier entry, although it occurred later in real time.

The absence of a dedicated tracing log is the primary finding for the required function. The substitute error log does collect evidence of connection-loss events, but it carries only a relative system tick and no actor identity. The system tick resets at a cold restart, so entries cannot be ordered across a power cycle. An entry recorded after the restart can carry a lower tick than an earlier entry, which confirms the reset. The log content also did not persist across the cold restart: the two entries from the first event were absent afterwards. This evidence-based observation directly demonstrates cold-start deletion in the substitute artefact and supports the conclusion that the retention capability is not demonstrated for RQ-012.


\subsection{Availability and Resilience}
\label{subsec:cat_availability} 
This category covers the availability aspect of RQ-007 and the resilience aspect of RQ-010, with the availability and domain-separation assets. The abstract catalogue contains the baseline resilience scenario, the network and layer-two flood scenario, the application-layer flood and domain-separation scenario, and the connection-hold self-denial-of-service scenario. The application-layer load test, the connection-hold procedure, and the layer-two flood were attempted under \gls{acr:NV1}. The domain-separation confirmation depends on observation of the safety channel and was designed but not executed, because it requires the \gls{acr:NV2} position. The activities executed under \gls{acr:NV1} are reported below.


\paragraph{Availability under Connection Load}
\label{subsec:res_availability}
This activity examined the availability of the identification interface under connection load. It relates to CM-009, to requirements RQ-007 and RQ-010, and to assets H and F. It targets the \gls{acr:IoT}-Core \gls{acr:HTTP} service and was run under \gls{acr:NV1}. The commands are listed in Appendix \ref{app:cmd_availability}.

A load test with Siege was run against the \gls{acr:IoT}-Core \gls{acr:HTTP} service at 50 and at 3 concurrent connections. The two load levels reflect the framework default and a near-envelope stress case for the \gls{acr:SRIO}, whose requirement documentation permits at most two concurrent connections. The console output and the Siege summary logs are stored as evidence.

The acceptance criterion is bounded by the specified operating envelope: availability is assessed only for the number of concurrent connections that the product specification permits. Both the 3-connection and 50-connection runs are therefore recorded as outside-envelope observations. Neither can, by itself, support a Pass or Fail verdict for the specified operating condition.

The stored summaries show that the service responded to a large number of requests with \gls{acr:HTTP} status 200 and that a substantial number of connections failed in the same runs. The six 50-connection runs recorded 1,273 successful transactions and 407 failed connections in total. One run completed no transaction and failed all 100 connection attempts after 26.58 seconds. The 3-connection run recorded 332 successful transactions and no failed connections. These results document behaviour under the tested loads, but the tested loads exceeded the specified maximum and therefore do not establish availability within the product's declared operating envelope.

The connection-hold procedure was also performed against the identification interface. Two long-lived \gls{acr:HTTP} connections were opened and are recorded as active in the job listing. While both connections were held, a third request to the identification endpoint did not receive a response within the five-second timeout. After one held connection was released, the same request again returned \gls{acr:HTTP} status 200. The evidence therefore demonstrates the intended connection-hold condition and recovery of the endpoint after a slot became available. The third request was made while the two permitted concurrent connections were already occupied and is consequently an outside-envelope overload observation, not a direct test of the specified operating condition. The activity remains \textbf{Not Assessed} for the requirement because it exercised only one endpoint and does not establish availability of all required identification information under every permitted operating condition.

The load results and the connection-hold attempt do not, by themselves, support a Pass or Fail verdict. A statement about domain separation or about the safety function is not possible under \gls{acr:NV1}, because the safety channel was not observed.

\paragraph{Layer-Two and Network Flooding}
\label{subsec:res_l2flood}
This activity examined the resilience of the device against a link-layer and network flood. It relates to requirement RQ-010 and to assets H and F. It targets the shared Ethernet segment and was run under \gls{acr:NV1}. The commands are listed in Appendix \ref{app:cmd_flood}.

The planned packet and \gls{acr:MAC}-address flooding activity is not counted as an evidence-backed execution. No complete terminal capture, packet capture, or independent record of the reported device indication is available. The activity is therefore recorded as blocked for assessment purposes and does not support a requirement verdict. To observe the effect of such a disturbance on the cyclic safety communication, the passive-observation variant NV2 would be required.

The robustness aspect of RQ-010 under malformed input is examined with the same availability asset. The abstract catalogue contains the protocol-fuzzing scenario for the fieldbus-discovery, safety-frame, and web-interface inputs. The fieldbus-discovery and safety-frame fuzzing require the \gls{acr:NV2} or \gls{acr:NV3} position and were designed but not executed. The web-interface fuzzing executed under NV1 is reported below.


\subsection{Robustness}\label{subsec:cat_robustness} This category covers the robustness aspect of RQ-010 under malformed input, with the availability asset H. The abstract catalogue contains the protocol-fuzzing scenario for the fieldbus-discovery, safety-frame, and web-interface inputs. The fieldbus-discovery and safety-frame fuzzing require the \gls{acr:NV2} or \gls{acr:NV3} position and were designed but not executed. The web-interface fuzzing executed under \gls{acr:NV1} is reported below.


\paragraph{Robustness under Malformed Input}
\label{subsec:res_robustness}
This activity examined the robustness of the \gls{acr:HTTP} service against malformed input. It relates to CM-010, to requirement RQ-010, and to asset H. It targets the \gls{acr:IoT}-Core \gls{acr:HTTP} service and was run under \gls{acr:NV1}. The commands are listed in Appendix \ref{app:cmd_robustness}.

A fuzzing test was run with boofuzz against the \gls{acr:HTTP} service on port 80. The saved script fuzzes the request line of an \gls{acr:HTTP} GET request. The boofuzz session databases and the run log are stored as evidence. The log records successful processing of the initial test cases, followed by a failed reconnection after a subsequent malformed request. The \gls{acr:HTTP} service became reachable again approximately 41 seconds later. The triggering input and the internal cause of the interruption were not established.

The stored evidence therefore shows a temporary disruption of the \gls{acr:HTTP} service on the non-safety path under malformed input, followed by recovery. The disruption requires follow-up by development. Its causation, reproducibility, and effect on the safety function are not established, because the safety channel was not monitored. The activity therefore supports an observation but not a complete requirement verdict for RQ-010; the requirement remains \textbf{Not Assessed} under the conservative aggregation rule.


\subsection{Authorisation and access control}
\label{subsec:cat_access}
This category covers the identification-and-authentication and the authorisation aspects of \gls{acr:FR} 1 and \gls{acr:FR} 2, addressing the modification-prevention and access-restriction requirements RQ-011 and RQ-014, together with the configuration-modification evidence of RQ-009, with the operating-mode, configuration, and audit-data assets. The abstract catalogue contains the read-only enforcement scenario, the unauthenticated configuration-protocol reset, the fieldbus parameter-write scenario, the firmware-update mode gate, the physical-barrier scenario, the unauthenticated tracing-log read, and the access-control and deletion-control probes. The device exposes a single anonymous read-only role and no authentication endpoint; the authentication dimension of this category therefore reduces, for this device, to the empirical finding that no identity model exists to authenticate against. The configuration-protocol and physical-barrier scenarios were designed but not executed. The read-only write probe and the access-control probe executed under \gls{acr:NV1} are reported below.

\paragraph{Access Control and Configuration Interfaces}
\label{subsec:res_access_control}
This activity examined the access-control model of the non-safety interface. It relates to CM-003 and CM-004, to requirements RQ-011, RQ-014, and RQ-009, and to assets B and G. It targets the \gls{acr:IoT}-Core interface and the configuration protocol and was run under \gls{acr:NV1}. The commands are listed in Appendix \ref{app:cmd_access}.

The \gls{acr:IoT}-Core interface provided anonymous read access without a login. CM-003 was not applicable because no authentication endpoint was exposed. No Hydra execution is claimed. An active web scan was run with \gls{acr:OWASP} ZAP through a local proxy, and the report was saved. A comparison of multiple user roles was not applicable, because the interface exposes a single anonymous read-only role.

A write probe was run against three \gls{acr:IoT}-Core nodes. A POST request to the safety-address node, to a digital-output node, and to the application-tag node was rejected in each case, and the responses are stored as evidence. The device answered with a bad-request response rather than an explicit authorisation denial. A bad-request response shows that the write was not accepted, but it does not by itself prove an access-control decision, because it can also result from the request format. For the application-tag node, which the device documentation lists as writable, an explicit authorisation denial would have been the expected response.

The anonymous read-only access is an observation about the \gls{acr:IoT}-Core interface. It does not, on its own, decide RQ-011 or RQ-014. The three submitted write requests did not result in an accepted modification, but the returned bad-request responses do not distinguish input validation from an authorisation decision. The activity therefore does not demonstrate an access-control mechanism. It does not cover the full write surface, the configuration-protocol behaviour, or the factory-reset path, which were not exercised. The tracing-log access control of RQ-014 is not evidenced beyond the single anonymous read-only role. The requirements therefore remain \textbf{Not Assessed} under the conservative aggregation rule.

\subsection{Open Test Areas}
\label{subsec:cat_config_firmware}
This subsection collects the firmware and configuration-integrity category together with the test areas that were specified in the methodology but not executed in the practical phase, so that they are not misread as results. They comprise the firmware and configuration-integrity tests, the tests that require the \gls{acr:NV2} or \gls{acr:NV3} position, and the company tests that were not applicable or blocked.

The firmware and configuration-integrity tests cover the integrity and authenticity of safety-relevant software and data under RQ-005, the parameter-integrity aspect of RQ-002, and the physical tamper resistance of the address and mode interface. They target the configuration, firmware, and operating-mode assets. The abstract catalogue contains the parameter-integrity rejection and forgery scenarios, the unsigned and modified firmware installation scenario, the update-package static-analysis scenario, the analytical safety-self-test scenario, and the physical tamper-resistance scenario. These scenarios require the engineering tool, a spare unit for a risk-bearing firmware test, the \gls{acr:NV3} in-path position, or physical access to the sealed interface. They were designed but not executed in the practical phase, so no requirement in this category is assessed from stored device evidence. The parameter-integrity, firmware-authenticity, and physical-tamper gaps identified during the theoretical validation remain open items for a future test campaign.

\label{subsec:res_na_blocked}
Three company tests could not be applied to the device, or were blocked by the environment. CM-002 requires a second zone and a conduit model, which the flat laboratory segment does not provide. CM-006 requires an operating-system shell, which the embedded device does not expose. CM-008 could not be operated, because the OpenVAS/GVM scanner could not be made operational within the laboratory setup and the available time. These outcomes are recorded as applicability or environment findings. They are meaningful methodological results, because they show why an applicability assessment is required before execution. They do not indicate a product defect, and they do not support a positive or a negative requirement verdict.

The following evidence items remain open: confirmation of the probable Spring Boot Actuator false positive by a direct request; examination of the \gls{acr:SNMP} service on port 161; and verification of whether aggressive Nmap timing destabilises device responses. The connection-hold procedure is no longer open as an evidence item: the new records demonstrate two occupied connections, timeout of a third request, and successful recovery after release. These observations remain bounded by the tested endpoint and operating conditions.




\section{Requirement Assessment}
\label{sec:requirement_verification}
The purpose of this section is to show that the methodology can assign a traceable verdict to each derived component-level requirement from the collected evidence. The per-requirement results are the output of applying the method to one device under the NV1 position. They are not a conformance audit of the product, and no legal conformity is claimed. A Fail verdict denotes an evidenced absence of component capability against a derived requirement, not a declaration of legal non-conformity of the product. Under the allocation defined in Section \ref{subsec:regulatory_applicability}, the corresponding machine-level obligation may instead be satisfied by a compensating countermeasure at system or zone level, which the machine manufacturer or integrator must document. Section \ref{subsec:req_verification_table} relates the executed activities of Section \ref{sec:executed_tests} to the fourteen derived requirements and assigns a verdict according to the scale of Section \ref{subsec:verdicts}. Section \ref{subsec:metrics_application} summarises the resulting coverage of the practical phase as direct counts.


\subsection{Verification Status per Requirement}
\label{subsec:req_verification_table} 
All executed activities reported here used the network variant \gls{acr:NV1}. A verdict is assigned only where a stored artefact supports it. The execution status of a test and its verdict are distinct. An executed test can still leave the requirement Not Assessed when the complete acceptance condition is not covered. A definitive Fail can result from an evidenced device deviation even though the test itself succeeded. Where the evidence documents such a deviation, the requirement is assessed as Fail because the component capability is not demonstrated against the derived requirement. Table \ref{tab:req_verification} summarises the current status per requirement. No requirement is verified as fully met. No legal conformity is claimed. Incomplete or ambiguous evidence is reported as \textbf{Not Assessed} under the conservative aggregation rule of Section \ref{subsec:requirement_assessment}.
%TODO Tabelle überprüfen
\begin{xltabular}{\textwidth}{l X l X} 
	\caption{Current requirement verification status under NV1} 
	\label{tab:req_verification} \\ 
	\toprule 
	\textbf{Req.} & \textbf{Supporting activity} & \textbf{Assessment} & \textbf{Limitation} \\ 
	\midrule 
	\endfirsthead 
	\caption[]{Current requirement verification status under NV1 (continued)} \\ 
	\toprule 
	\textbf{Req.} & \textbf{Supporting activity} & \textbf{Assessment} & \textbf{Limitation} \\ 
	\midrule 
	\endhead 
	\midrule 
	\multicolumn{4}{r}{\textit{Continued on next page}} \\ 
	\endfoot 
	\bottomrule 
	\endlastfoot 
	RQ-001 & attack-surface enumeration (Section \ref{subsec:cat_identification}) & Not Assessed & non-safety path only; no safety-channel or complete domain-separation evidence \\ 
	RQ-002 & none executed & Not Assessed & requires the engineering tool or the relevant NV2/NV3 scenario \\ 
	RQ-003 & intervention evidence and time attribution (Section \ref{subsec:cat_logging}) & Fail & no dedicated tracing log; no absolute timestamp or actor identity; compensating control at system level required \\ 
	RQ-004 & identification inventory (Section \ref{subsec:cat_identification}) & Not Assessed & product-level identification was observed, but complete software and data coverage was not established; cleartext is non-scoring evidence \\ 
	RQ-005 & none executed & Not Assessed & firmware and integrity tests not executed \\ 
	RQ-006 & installed-software inventory (Section \ref{subsec:cat_identification}) & Not Assessed & product granularity only; complete component inventory was not established \\ 
	RQ-007 & identification availability; connection-hold and load tests (Sections \ref{subsec:cat_identification}, \ref{subsec:cat_availability}) & Not Assessed & availability and recovery were observed at the tested endpoint, but the complete specified operating envelope was not assessed \\
	RQ-008 & software-event evidence baseline (Section \ref{subsec:cat_logging}) & Not Assessed & baseline evidence exists, but no software-installation event was induced \\ 
	RQ-009 & configuration evidence baseline (Section \ref{subsec:cat_logging}) & Not Assessed & baseline evidence exists, but no configuration change was induced \\ 
	RQ-010 & HTTP fuzzing; load test (Sections \ref{subsec:cat_robustness}, \ref{subsec:cat_availability}) & Not Assessed & safety channel was not observed and the layer-two activity was not evidence-backed \\ 
	RQ-011 & write probe on three nodes (Section \ref{subsec:cat_access}) & Not Assessed & bad-request responses do not prove an access-control decision \\ 
	RQ-012 & cold-start deletion (Section \ref{subsec:cat_logging}) & Fail & substitute error log cleared by a cold restart; compensating control at system level required \\ 
	RQ-013 & none executed & Not Assessed & version history not exposed; current-version read does not evidence retention \\ 
	RQ-014 & access-control probe (Section \ref{subsec:cat_access}) & Not Assessed & single anonymous role; tracing-log restriction and unauthorised observation were not fully assessed \\ 
\end{xltabular} 
	
The strongest evidence concerns the intervention log. The device records connection-loss events, but no dedicated tracing log was exposed through the examined interface; the substitute error-log entries carry only a relative system tick, no actor identity, and the log content is cleared by a cold restart. The evidence therefore supports the conclusion that the required tracing capability is not demonstrated and that the substitute artefact has the documented limitations behind RQ-003 and RQ-012. The identification tests demonstrate product-level information retrieval, but they do not establish complete software or data coverage. Cleartext HTTP is recorded as a defence-in-depth observation and does not itself verify RQ-004 or RQ-014. RQ-013 is Not Assessed because a current-version read does not evidence retention of uploaded software versions. The robustness and load observations do not provide a complete RQ-010 assessment because the safety channel was not observed. The write probe shows only that three specific write requests were not accepted; because the device answered with a bad-request response, it does not establish an access-control decision. The remaining requirements are Not Assessed because their complete acceptance conditions were not covered under NV1.


\subsection{Coverage of the Practical Phase}
\label{subsec:metrics_application} 
The assessment model of Section \ref{sec:evaluation_metrics} is qualitative, so the results are summarised as direct counts rather than as ratios. Fourteen requirements were derived from the Regulation. Of the ten company tests, five were applicable and executed, three were not applicable to the device or topology, and one was blocked in the laboratory; one further procedure was considered applicable but the original test case was not usable as a direct security verdict because it primarily produced a transport-level observation. The exact number of executed test activities is determined from the activity-and-evidence register: an activity is counted once only when a stored artefact supports its interpretation. Unsupported observations and planned procedures are excluded. Two requirements, RQ-003 and RQ-012, are assessed as Fail on the basis of evidenced component-capability deviations. The remaining requirements are Not Assessed because their complete acceptance conditions were not covered.
The evidence base includes raw Nmap and Nikto output, SSLscan output, inventory reads, error-log and system-tick reads, Siege summaries, connection-hold output, write-probe responses, boofuzz session data, and the \gls{acr:OWASP} ZAP report. These artefacts support reproducible observations, but not necessarily complete requirement verification. The limited number of definitive verdicts reflects the bounded scope of the practical phase. The assessment can be extended when the outstanding \gls{acr:NV2}, \gls{acr:NV3}, firmware, version-history, and engineering-tool activities are executed and recorded.






\section{Discussion of Results}
\label{sec:discussion} 
Section \ref{sec:requirement_verification} assigned a verdict to each of the fourteen derived requirements. This section steps back from the individual verdicts and interprets the resulting pattern as a whole. Section \ref{subsec:discussion_synthesis} discusses why the practical phase yields predominantly Not Assessed verdicts, what the two Fail verdicts signify, and what role the unrealised network variants \gls{acr:NV2} and \gls{acr:NV3} play in this pattern. Section \ref{subsec:discussion_rq} relates these findings back to the three research questions defined in Section \ref{subsec:research_questions}.

\subsection{Synthesis of Findings Across Test Categories}
\label{subsec:discussion_synthesis} 
Table \ref{tab:req_verification} shows a pattern rather than an isolated result: two requirements are assessed as Fail, and the remaining twelve are Not Assessed. Under the conservative aggregation rule of Section \ref{subsec:requirement_assessment}, a Not Assessed verdict does not indicate that a security function is absent. It indicates that the executed evidence does not cover the complete acceptance condition of the requirement. The predominance of this verdict is therefore not primarily a finding about the device, but a direct consequence of the scope decisions documented in Section \ref{subsec:practical_scope}: the practical phase realised only the direct-access variant NV1, while a majority of the abstract test scenarios in Section \ref{sec:executed_tests} that could have produced a Pass or a further Fail depend on \gls{acr:NV2}, \gls{acr:NV3}, an induced software or configuration change, or the engineering tool. This is visible category by category. In identification and asset discovery (Section \ref{subsec:cat_identification}), product-level identification was demonstrated, but a complete \gls{acr:SBOM} or data-signature check was out of \gls{acr:NV1} reach. In availability and resilience and in robustness (Sections \ref{subsec:cat_availability}, \ref{subsec:cat_robustness}), the load test and the fuzzing activity produced observations, but neither could be related to the cyclic safety channel, which only \gls{acr:NV2} exposes. In authentication and access control (Section \ref{subsec:cat_access}), the write probe demonstrated that three specific requests were rejected, but a bad-request response does not distinguish input validation from an authorisation decision, so the requirement remains open rather than confirmed.

The two Fail verdicts, RQ-003 and RQ-012, stand apart from this pattern because they were fully assessable within \gls{acr:NV1} alone. Both concern the same substitute artefact, the \gls{acr:IoT}-Core error log, and both are demonstrated by the same evidence sequence in Table \ref{tab:rq003_evidence}: the log carries only a relative system tick and no actor identity, and its content is deleted by a cold restart. This is the strongest finding of the practical phase, precisely because it did not require passive observation of the safety channel or an in-path position to demonstrate. It illustrates a general property of the methodology: a requirement can be conclusively assessed under a restricted network position when the demonstrating property lies entirely on the non-safety, \gls{acr:IP}-reachable path. Conversely, requirements whose acceptance condition touches the cyclic PROFINET/PROFIsafe communication, such as the domain-separation aspect of RQ-001 or the resilience aspect of RQ-010, cannot be conclusively assessed under \gls{acr:NV1} regardless of how many \gls{acr:NV1}-compatible tests are executed against them, because the observation point itself is insufficient for the claim.

The unrealised \gls{acr:NV2} and \gls{acr:NV3} variants therefore do not merely reduce the number of executed test cases. They remove the observation points required for entire classes of claims, specifically those concerning the safety-relevant process data and in-path manipulation. This distinguishes two different reasons for a Not Assessed verdict that the aggregation rule of Section \ref{subsec:requirement_assessment} treats alike but that carry different implications for future work: a requirement can be open because the induced event was never triggered under an otherwise sufficient access position (for example RQ-008 and RQ-009, where only baseline evidence was collected), or it can be open because the required access position itself was never established (for example the domain-separation aspect of RQ-001 and RQ-010). The first class can be closed by extending the test execution within the existing \gls{acr:NV1} topology. The second class requires the topology itself to be extended, specifically a mirror or \gls{acr:SPAN} port for \gls{acr:NV2} or an inline bridge for \gls{acr:NV3}, as noted in Section \ref{subsec:practical_scope}.


\subsection{Answering the Research Questions}
\label{subsec:discussion_rq} 
\textbf{RQ1: How can the machine-level cybersecurity obligations of Regulation (EU) 2023/1230 be translated into verifiable, derived component-level security tests for a safety-related industrial gateway?} The instantiation in Section \ref{subsec:req_asset_interface} demonstrates that the derivation chain defined in Section \ref{subsec:derivation_mapping}, from the Regulation through prEN 50742 to \gls{acr:IEC} 62443-4-2, and further through the asset and threat mapping of Section \ref{subsec:asset_threat_mapping}, produces requirements that are concrete enough to bind to a specific interface, asset, and network access variant on a real device (Table \ref{tab:req_asset_interface}). The subsequent execution in Section \ref{sec:executed_tests} confirms that each requirement can be operationalised into at least one test activity, even where the activity ultimately yields a Not Assessed verdict rather than a Pass. The two Fail verdicts, RQ-003 and RQ-012, show that the chain is capable of producing a conclusive negative verdict when the acceptance condition lies within the realised access position. RQ1 is therefore answered affirmatively within the scope of NV1: the translation from a machine-level obligation to a verifiable, evidence-based component test is demonstrated end to end, though the demonstration is fully conclusive only for the subset of requirements whose acceptance condition does not depend on \gls{acr:NV2} or \gls{acr:NV3}. 

\textbf{RQ2: Which tests of the existing company framework are applicable to an industrial gateway, and how can non-applicable tests be identified before execution?} Table \ref{tab:cm_applicability} answers this question directly: of the ten company tests, five were applicable and executed (CM-001, CM-004, CM-007, CM-009, CM-010), three were not applicable to the device or the topology (CM-002, CM-003, CM-006), one original procedure was not applicable while the tool itself produced a usable interface observation (CM-005), and one was blocked by the laboratory environment (CM-008). The applicability assessment of Section \ref{subsec:cm_applicability} identified all three not-applicable cases before execution, on the basis of an explicit tool-to-target check: CM-002 assumes a second network zone that the flat topology does not provide, CM-003 assumes an authentication endpoint that the \gls{acr:SUT} does not expose, and CM-006 assumes an operating-system shell that the embedded device does not expose. RQ2 is therefore answered by the methodology itself: a systematic applicability assessment, performed prior to execution and grounded in the assumed interface of each tool rather than in the tool's general purpose, correctly separates applicable tests from tests whose assumed target interface is known in advance to be absent. This distinction is what prevents a not-applicable outcome from being misread as a passed or failed security test. 

\textbf{RQ3: What can direct-access testing at the physical \gls{acr:IT}/\gls{acr:OT} interface verify about the effectiveness of the implemented security functions, and where are its boundaries?} The synthesis in Section \ref{subsec:discussion_synthesis} bounds this answer precisely. Direct-access testing under \gls{acr:NV1} can conclusively verify requirements whose acceptance condition is exercised entirely through the \gls{acr:IP}-reachable non-safety path, as demonstrated by the Fail verdicts for RQ-003 and RQ-012. It can also produce reproducible interface-level observations, such as the absence of a \gls{acr:TLS} endpoint (Section \ref{subsec:cat_communication}) or the absence of an authentication endpoint (Section \ref{subsec:cat_access}), that constrain but do not by themselves decide a requirement. Its boundary is equally precise: \gls{acr:NV1} cannot observe the cyclic PROFINET/PROFIsafe communication between controller and \gls{acr:SUT}, so it cannot support any claim about the integrity of safety-relevant process data in transit, about in-path manipulation, or about the effect of a network disturbance on the safety function itself. Every requirement in Table \ref{tab:req_verification} whose limitation column names the safety channel, the \gls{acr:NV2} position, or the \gls{acr:NV3} position marks exactly this boundary. RQ3 is therefore answered as follows: direct-access testing is sufficient to verify security functions that are fully exposed on the \gls{acr:IT}-side interface, but it structurally cannot verify or refute claims about the \gls{acr:OT}-side safety communication, regardless of the number or the depth of the tests executed at the \gls{acr:IT}/\gls{acr:OT} interface.


\section{Evaluation of the Methodology}\label{sec:objective_validation} This section closes the chapter by evaluating the methodology itself, as distinct from the requirement verdicts of Section \ref{sec:requirement_verification} and their interpretation in Section \ref{sec:discussion}. Section \ref{subsec:objective_validation_sub} validates the methodology against the thesis objective, following the evaluation activity of the Design Science Research process. Section \ref{subsec:threats_to_validity} examines the reliability of the findings and the threats to construct, internal, and external validity that bound them.


\subsection{Validation against the Thesis Objective}
\label{subsec:objective_validation_sub} 

This section validates the developed methodology against the thesis objective. The evaluation follows the evaluation activity of the Design Science Research process \cite[p. 46]{Peffers}. The thesis objective is to develop and validate a systematic methodology that supports the verification of security functions for industrial gateways at the physical \gls{acr:IT}/\gls{acr:OT} interface. The methodology was applied end to end on a real device. It derived requirements from the regulation, connected them to assets and interfaces, identified applicable and non-applicable tests before execution, adapted the generic company procedures, collected evidence, and related observations to technical requirements. Table \ref{tab:goal_validation} maps each objective component to the methodology element that realises it, to the practical evidence, and to the validation boundary. 

\begin{table}[H] 
	\centering 
	\small 
	\caption{Validation of the methodology against the thesis objective} \label{tab:goal_validation} 
	\begin{tabularx}{\textwidth}{X X l} 
		\toprule 
		\textbf{Objective component} & \textbf{Methodology element and evidence} & \textbf{Status} \\ 
		\midrule 
		Derive requirements from the regulation & Machinery Regulation to prEN 50742 to \gls{acr:IEC} 62443-4-2 derivation; 14 requirements (Table \ref{tab:mvo-requirements}, Appendix \ref{app:requirement_derivation}) & demonstrated \\ 
		Connect requirements to assets and interfaces & asset mapping and asset table; requirement-to-asset-to-interface mapping (Table \ref{tab:req_asset_interface}) & demonstrated \\ 
		Assess security testing approaches & structured comparison and selected combination (Section \ref{subsec:approach_evaluation}) & demonstrated \\ 
		Evaluate the company framework & applicability assessment; Table \ref{tab:cm_applicability} & demonstrated \\
		Identify tool-to-target mismatches & CM-003 and CM-005 not applicable to the original procedures; CM-006 not applicable; CM-008 blocked & demonstrated \\ 
		Define requirement-oriented test cases & two-level catalogue; executed subset (Section \ref{sec:executed_tests}) & demonstrated within the defined scope \\ 
		Define test-environment variants & \gls{acr:NV1}, \gls{acr:NV2}, \gls{acr:NV3} model; \gls{acr:NV1} realised & demonstrated for \gls{acr:NV1} \\ Apply the method to a real object & \gls{acr:SRIO} instantiation; stored evidence & demonstrated \\ 
		Distinguish applicable, not applicable, blocked, and open tests & separate execution-status and requirement-verdict model; Tables \ref{tab:cm_applicability} and \ref{tab:req_verification} & demonstrated \\ 
		Support integration into future releases & lifecycle and governance model (Appendix \ref{app:SecureDevelopmentLifecycle}); reusable schema & specified by design \\
		\bottomrule 
	 \end{tabularx} 
 \end{table} 
 
The empirical validation covered the central methodology steps under \gls{acr:NV1}. The requirement derivation, the asset mapping, the applicability assessment, the test selection, the execution, and the evidence-based assessment were all exercised on a real device. How far this validation extends the specific capabilities of the methodology, expressed through the three research questions, is discussed in Section \ref{subsec:discussion_rq}. The tests that require passive observation or inline manipulation remain extensions of the validated core process. The methodology is therefore validated within a bounded scope, defined by the \gls{acr:NV2} and \gls{acr:NV3} variants and by the requirements whose evidence is still open. A claim of full validation is not made. The validation concerns the methodology rather than the product, and its boundary is the \gls{acr:NV1} position together with the requirements whose evidence remains open. These limits are examined further in Section \ref{subsec:threats_to_validity}.

\subsection{Reliability and Threats to Validity}
\label{subsec:threats_to_validity} 
The results are subject to threats to construct validity, internal validity, external validity and reliability. 

\textbf{Construct validity.} The fourteen requirements are an interpretation of machine-level obligations and were derived by a single author. They are mapped through the draft standard prEN 50742 to \gls{acr:IEC} 62443-4-2, so the interpretation may change if the final standard text differs. The mapping is therefore itself a threat to validity. Several practical observations are also indirect: the Nikto Actuator finding remains unconfirmed, a bad-request response does not distinguish input validation from authorisation, and the SSLscan result establishes the absence of a \gls{acr:TLS} endpoint but not a complete confidentiality assessment. 

\textbf{Internal validity.} The fuzzing-induced \gls{acr:HTTP} disruption was not reproduced, and its causation cannot be attributed to a specific input. The planned layer-two flooding activity lacks a complete terminal capture and is not treated as executed evidence. In addition, hping3 and macof would target the switch and the shared segment as well as the device, so any uncorroborated indication could not be attributed to a device-side security response. 

\textbf{Safety-priority constraint.} Section \ref{subsec:constraints_scope} requires that testing must not impair the safety function. Under \gls{acr:NV1}, however, the cyclic safety channel was not observable. The constraint was therefore maintained procedurally through the isolated laboratory and the exclusion of production systems rather than demonstrated by measurement. 

\textbf{Connection-hold result.} The new evidence establishes that two \gls{acr:HTTP} connections were occupied, that a third request timed out while both slots were held, and that the endpoint responded successfully after a slot was released. This confirms the connection-hold procedure and recovery behaviour, but the third request was outside the specified two-connection operating envelope and only one identification endpoint was exercised. The requirement therefore remains \textbf{Not Assessed}, because the evidence does not cover the complete acceptance condition.

\textbf{External validity.} The evaluation covers one device, one firmware revision, one flat laboratory topology without \gls{acr:VLAN} separation, and only the \gls{acr:NV1} network position. The results do not generalise to other gateways or to the cyclic safety path. \textbf{Reliability.} The work was performed by a single tester and depends on specific tool versions and scan timing. The Nmap timing sensitivity is itself evidence that execution conditions affect observations. The stored tool outputs, commands in Appendix \ref{app:tool_commands}, and defined acceptance criteria make the raw evidence reproducible, although the interpretation is not fully reproducible where causation or setup state is uncertain. These limitations bound the claims of Section \ref{subsec:objective_validation_sub} to a bounded validation of the methodology.
